%% OPEN additions retained after exact Formal Practice restoration.
% Necessity: the exact inherited chapter does not contain the worked
% two-element exercise solution, automation/encoding/mechanized-result methods,
% writing/production methods, per-result record, or release-gate apparatus.
% Definition: this file is exactly the three complete marked historical slices
% below with no source-token rewrite or status promotion.
% Standing: it is an administrative OPEN body supplement, not inherited source,
% an admission, an inference rule, or a completed release review.
% Nonredundancy: the five exact wrappers provide every inherited method, example,
% warning, failure table, and exercise; these slices retain only a separate
% worked solution and later audit/release apparatus from the historical chapter.
% Use: parts/p12.tex inputs this file once after SRC-14231-14270; route, slice,
% label, consumer, and publication tests fail closed on any drift.
\noindent\textit{Reader note.} The following worked exercise solution and later
audit/release methods are separately attributed OPEN additions to the exact
inherited chapter above; they neither replace nor alter its methods or claims.
The editorial audit methods also do not replace either later exact chapter:
Writing and Typesetting Formal Logic or Research Workflow and Reproducibility.

% BEGIN HISTORICAL SLICE 50-59
\begin{hbexample}{def:p12-two-element-countermodel}{Why one element fails and two suffice}
The premise \(\forall x(P(x)\lor Q(x))\) does not semantically entail
\((\forall xP(x))\lor(\forall xQ(x))\).  On a one-element domain with element
\(a\), the premise makes \(P(a)\) or \(Q(a)\) true, which makes one disjunct of
the conclusion true; hence no one-element countermodel exists.  On
\(\{0,1\}\), interpret \(P\) as \(\{0\}\) and \(Q\) as \(\{1\}\).  Every
element satisfies \(P(x)\lor Q(x)\), but neither predicate holds of every
element, so the conclusion fails.
\end{hbexample}

% END HISTORICAL SLICE 50-59

% BEGIN HISTORICAL SLICE 73-192
\section{Automation and encoding audits}

\begin{hbmethod}{def:p12-automated-result-audit}{Automated-result audit}
\index{automated-result audit}
Record the exact input artifact or embedded input text and its content hash,
the input language and version, intended semantics, every declaration and
nonlogical axiom, each transformation and whether it preserves
formula equivalence or only equisatisfiability, the calculus and search
program, tool build and options, exact invocation command and working
directory, and, for every generated example, a fixed random seed.  Record the
expected response or exit status, result token, resource limit,
model or proof certificate, and the independent checker's exact version,
invocation or checking steps, expected result, and observed result.  Also
record the trusted base and the result of reproduction in a declared clean
environment.  This record is necessary
because a timeout, \texttt{unknown}, unchecked model, and checked refutation
have different evidential standings.  It is used for SMT-LIB, TPTP, and SAT
reports in Chapters~\ref{sec:p12-sat-and-certificates} and
\ref{sec:p12-smt-and-exchange-formats}.
\end{hbmethod}

\begin{hbmethod}{def:p12-encoding-audit}{Encoding audit}
\index{encoding audit}
State the source and target syntactic categories as typed maps; state whether
the map is injective, surjective, compositional, size-bounded, or none of these;
and state the exact preservation/reflection claim needed.  Audit the
representations of the signature, terms, formulas, and binding; the definition
and correctness obligations for capture-avoiding substitution; and the exact
policy for alpha-equivalence or quotienting.  A \emph{de Bruijn index} is a
natural-number code for a bound occurrence given by its binder distance under
an explicitly declared zero- or one-based convention.  It stands as a
nameless binding representation, not the independent-small-checker question
recorded in Table~\ref{tab:p12-assistant-landscape}, and it is not merely
an alpha-equivalence quotient.  When indices are used, the audit must record
the base convention and the capture-avoiding shifting and substitution
operations.  For a deep
embedding, audit constructors, structures and satisfaction or derivations,
and the interpretation; for a shallow embedding, audit which host equalities
or operations identify object expressions.  State the correspondence theorem
in both directions.  For a reflected procedure, record the internal soundness
theorem from Definition~\ref{def:p12-computational-reflection}; also record all
axioms and unsafe, opaque, native, or externally trusted extensions through
the mechanized-result report.  The output is a list of correspondence
obligations, not an adequacy theorem; the distinction prevents the two
strategies in Definition~\ref{def:p12-deep-shallow-embedding} and
Example~\ref{def:p12-deep-shallow-conjunction} from being declared faithful by
notation alone.
\index{de Bruijn index}
\end{hbmethod}

\begin{hbmethod}{def:p12-mechanized-result-report}{Mechanized-result report}
\index{mechanized-result report}
Record the informal claim, exact formal theorem statement, correspondence
between them, imported definitions and libraries, axiom surface, assistant and
version, enabled options, whether the retained artifact is a proof term,
another proof object, or only a proof script, its kernel-checking path, external
oracles, compilation result, and independently reproduced checks.  When
computational reflection is used, record its internal representation and
interpretation, soundness theorem, and kernel-reduction, proof-producing, or
externally trusted evaluation path as fixed in
Definition~\ref{def:p12-computational-reflection}.  Report human mathematical
review separately.  This result-relative report is not a
product description and does not infer trust from the words “small kernel” or
“formally verified”; it supplies the fields used by
Tables~\ref{tab:p12-assistant-landscape} and
\ref{tab:p12-landmark-formalizations}.
\end{hbmethod}

\section{Writing and production audits}

\begin{hbmethod}{def:p12-formal-writing}{Formal-definition writing method}
\index{formal definition!writing method}
Replace every vague technical noun by its typed category, quantified
parameters, exact defining clause, and scope.  Replace “clearly” or “obviously”
by the definition, calculation, or cited result that licenses the step.
Separate existence, uniqueness, equivalence, and one-way implication.  At the
first occurrence of a new name, record its necessity, definition or axiom,
standing, nonredundancy, and later use.  The output is a candidate definition
or statement for review, not evidence for its truth.
\end{hbmethod}

\begin{hbmethod}{def:p12-notation-gate}{Notation-registry gate}
\index{notation registry}
For every mathematical glyph, record its role, type, scope, precedence, first
use, and collisions.  Reuse a registered role when its
meaning agrees.  When a named logic, model, or method has a fixed source
presentation, retain its published glyph and register that locally typed role;
for notation introduced only by this handbook, add a scoped row only after
checking overlaps.  Argument shape may disambiguate declared uses
of \(\entails\), but a glyph may not change standing silently within one scope.
This method differs from visual proofreading: it governs semantic roles before
rendering.  The complete editorial registry is printed in Part~1, and the
inherited source presentation remains reader-facing in
Chapter~\ref{ch:registered-notation}.
\end{hbmethod}

\begin{hbmethod}{def:p12-source-render-audit}{Source, render, and cross-reference audit}
\index{source audit}\index{rendered-page audit}\index{cross-reference audit}
First compare the exact statement and status against its controlling source.
Then compile until labels and citations resolve, run the static status/sheet
checks, inspect the rendered theorem, glyphs, tables, links, and page breaks,
and extract text to detect missing or corrupted content.  Finally follow every
body-to-sheet and body-to-annex route.  These are distinct checks: clean TeX
does not prove a theorem, a visible page does not establish source fidelity,
and text extraction does not replace visual inspection
\cite{AMSAuthorHandbook,AMSMathGuide,Unicode2025}.
\end{hbmethod}

\begin{hbmethod}{def:p12-counterexample-first-review}{Counterexample-first independent review}
\index{counterexample-first review}
Before reading the proposed proof, test empty premises, one- and two-element
structures, empty and universal relations, vacuous binding, inconsistent
theories, finite/infinite and countable/uncountable boundaries, and constructive
versus classical metalogic where relevant.  A second pass traces every
definition through the argument; a third compares the final statement with its
formal or computational encoding.  The method stands as an adversarial review
order, not as proof by finite testing, and its separate passes are recorded in
the release table.
\end{hbmethod}

% END HISTORICAL SLICE 73-192

% BEGIN HISTORICAL SLICE 219-348
\begin{hbmethod}{def:p12-result-status-record}{Per-result research status record}
\index{result status record}
A migrated or newly researched result requires a six-field record because the
generated audit register records only its label, title, evidence status, and
target; that register cannot state the exact claim, dependencies, formal
system, or verification actually performed.  The record is the following
metalinguistic research artifact, not an additional theorem status.
\begin{center}
\small
\begin{tabular}{@{}p{.20\textwidth}p{.70\textwidth}@{}}
\toprule
Field & Required content \\ \midrule
Identifier & The unique handbook result label or the exact open register row. \\
Exact statement & The complete quantified statement, including hypotheses,
active calculus or model class, and the object bearing the asserted property. \\
Dependencies & Every local result, registered sheet, axiom, or cited theorem
used to support the statement. \\
Status & Exactly one permitted evidence class for a result, or an explicit
open/owed-artifact disposition when no result is asserted. \\
System & The exact logic, calculus, model class, prover, or checker relative
to which the statement is made. \\
Verification & The separately reported checks actually performed, such as
source comparison, annex-proof review, kernel or certificate checking,
compilation, and rendered-page inspection. \\
\bottomrule
\end{tabular}
\end{center}
Its nonredundant role is to prevent a build-success or citation field from
standing in for mathematical and system-specific evidence.  The release audit
below requires this record for every migrated or researched result.
\end{hbmethod}

\section{Release gates}

\begin{landscape}
\thispagestyle{plain}
\begingroup\small
\begin{longtable}{@{}p{1.40in}p{2.35in}p{3.35in}@{}}
\caption{Formal-logical release audit.  Each gate is reported separately.}
\label{tab:p12-release-audit}\\
\toprule
Gate & Establishes & Does not establish \\ \midrule
\endfirsthead
\caption[]{Formal-logical release audit.  Each gate is reported separately (continued)}\\
\toprule
Gate & Establishes & Does not establish \\ \midrule
\endhead
\midrule\multicolumn{3}{r}{\emph{continued on next page}}\\
\endfoot
\endlastfoot
Definition and notation audit & First-use obligations, exact signatures and
grammars, typed roles, scoped glyphs, binding, and precedence declarations &
Existence or consistency of the defined objects \\ \addlinespace
Recursive-definition audit & A stated decreasing measure or other
well-foundedness argument for every recursive clause & Correctness of later
theorems using the recursively defined object \\ \addlinespace
Quotient-definition audit & Operations and predicates are independent of the
chosen representative of each equivalence class & Any additional universal
property or completeness claim \\ \addlinespace
Theorem-scope audit & Every quantifier, hypothesis, active calculus or model
class, and the exact object bearing the asserted property are explicit & Truth
of the resulting statement \\ \addlinespace
Induction and case audit & The induction object or measure and every base,
constructor, limit, and side-condition case actually required are present &
Correctness of imported lemmas not separately checked \\ \addlinespace
Binding/substitution audit & Quantifier-rule side conditions and
capture-avoiding substitution obligations are stated and discharged where
used & Semantic adequacy of the whole calculus \\ \addlinespace
Transformation/encoding audit & Method~\ref{def:p12-encoding-audit} records
typed source and target syntax, binding treatment, and whether each step
preserves formula equivalence, only equisatisfiability, or a separately stated
correspondence & Any stronger preservation theorem not recorded there \\
\addlinespace
Evidence-status audit & A sketch is not promoted to a proof; citations and
annex targets match; relative-consistency conclusions retain their base theory
and relativity & Correctness of a cited source or completion of an explicitly
open artifact \\ \addlinespace
Per-result status record & All six fields of
Method~\ref{def:p12-result-status-record} are present & Any verification not
listed in that record \\ \addlinespace
Status/static checks & Duplicate labels, macros, and glyph roles are rejected,
and label, status, sheet, stub, and register-shape invariants hold & Theorem
truth or completeness of an annex argument \\ \addlinespace
Source audit & Original mathematical claims use primary sources where
available, software claims use current authoritative documentation, and exact
editions or versions are recorded & Independent proof of the cited claims \\
\addlinespace
Repository-generation audit & Duplicate macros and glyph roles are resolved,
and generated symbol tables, audit tables, index, and bibliography agree with
the live sources & Mathematical correctness \\ \addlinespace
Automated-result audit & Method~\ref{def:p12-automated-result-audit} records
the exact input, tool/version/command, returned status, resource limit,
certificate or model, checker invocation and result, and trusted base & A
mathematical conclusion stronger than the checked artifact and status \\
\addlinespace
Mechanized-result audit & Method~\ref{def:p12-mechanized-result-report} records
the exact mechanized statement, representation/interpretation, correspondence,
imports, axioms or unsafe extensions, version, kernel, and checking result &
Fidelity to an informal theorem not established by the correspondence fields \\
\addlinespace
Clean compilation & TeX integration and resolved reference/bibliography paths &
Mathematical correctness \\ \addlinespace
Source-fidelity review & Agreement of migrated claims and evidence labels with
the controlling source & Correctness of genuinely new mathematics \\ \addlinespace
Mathematical review & Valid hypotheses and reasoning for the claims actually
checked & Clean typography or exhaustive correctness of unchecked claims \\
\addlinespace
Rendered-page inspection & Visible glyphs, clipping and overlap, line breaks,
tables, links, headings, headers, footers, page numbers, rule sheets, and page
integrity & Source semantics by itself \\ \addlinespace
Certificate or kernel check & Acceptance by the declared checker under the
reported axiom/import surface & Fidelity of the formal definitions to informal
intent or correctness of unchecked external code \\ \addlinespace
Automated-agent review & Findings under the reported agent, configuration,
scope, and artifacts & Human review or an assurance that no defect exists \\
\addlinespace
Reviewer-provenance audit & The identity, role, model or tool version,
configuration, scope, and exact artifact snapshot of each automated or human
review are recorded separately & Agreement among reviewers or correctness
beyond their declared scopes \\ \addlinespace
Human pedagogical review & Usability, navigation, and explanatory adequacy for
the reviewed audience & Exhaustive formal correctness \\ \bottomrule
\end{longtable}
\endgroup
\end{landscape}

The exact inherited notation chapter, rule sheets, assessment material,
exercise feedback, examples, exercises, and supplied solutions remain
reader-facing in their restored source order, including
Chapters~\ref{ch:registered-notation}, \ref{ch:rule-sheets}, and
\ref{ch:exercise-feedback}.  Cross-references to Part~1, the calculus registry,
Part~14, or the examples in this chapter are navigation only: they do not
replace or consolidate those source passages.
% END HISTORICAL SLICE 219-348
